% Atomic-response to finite-volume Bill transfer and Coffer price conversion.
\subsection{Atomic responses and the finite-volume root price}
\label{sec:ex-lattice-trace-coffer-supplier}

This section records the actual transfer used by the public trace supplier.
There are two distinct operations. An atomic response first becomes a
finite-volume response with a root-price estimate. The price estimate is
then converted into the two components of a Bill. We give the conversions
and the quantitative endpoint of the Coffer argument, while keeping the
upstream response, paid-price, and price-duality receipts explicit.

\subsubsection{From an atomic response to the Coffer input}

Let \(X\) be an order-zero coefficient at scale \(k\), \(g\) its local field,
and \(Z=\widehat\beta_k\). Write \(x_X\) for its coordinate, \(X^{\rm f}\) for its forced part, and \(D_X\) for its density. The atomic certificate supplies a nonnegative
remainder allowance \(s\), a response \(L\) with attainment cost at most
\(e\), and
the identities
\begin{align}
 \|x_X\|_{\rm coeff}+s&\leq q, \label{eq:ex-trace-atomd-charge}\\
 \operatorname{profile}(X^{\rm f})&=\operatorname{hook}(L),\\
 D_X(U)&=\operatorname{eval}(U)\,g(U).
 \label{eq:ex-trace-atomd-data}
\end{align}
The remainder bound is measured at scale \(k\), at order zero and with
tolerance \(\sqrt{\operatorname{forcedAllowance}_k}\). The statement that
the response \(L=(L_Y)_Y\), indexed by supports \(Y\), has attainment cost
at most \(e\) is a certificate, not a scalar equality: every component
\(L_Y\) is measurable and gauge invariant, vanishes at every configuration
that is healthy on \(Y\) at the threshold of the scale-\(k\) budget, and
takes the same value at two configurations whose links on \(Y\) agree and
whose threshold flags in \(Y\) agree; \(L_Y\neq0\) only if \(Y\) contains
the marked origin, and there are numbers \(a_Y\geq0\) with
\[
 |L_Y(U)|\leq a_Y\ \text{ for all }U,\qquad
 \sum_Y e^{(\alpha_k+2)\operatorname{cov}_k(Y)}a_Y\leq e,
\]
where \(\alpha_k\) is the support exponent and \(\operatorname{cov}_k\)
the cover number of that budget. Set
\(u=\operatorname{alto}_{k,d}\), the inverse field dilation, and \(G'=ug\).
Multiplying the density identity by \(u\) gives
\[
 uD_X(U)=\operatorname{eval}(U)G'(U).
\]
Thus the atomic response becomes the Coffer input, with unchanged format,
remainder and forced-profile data.

The selected atomic charge budget satisfies
\begin{equation}
 4(1+c_k)q\leq Z^3,
 \label{eq:ex-trace-coffer-budget}
\end{equation}
where \(c_k\geq1\) is the selected column allowance of the history at
scale \(k\), a real parameter of the stage data and not a cardinality.
Hence
\[
 4(1+c_k)(\|x_X\|_{\rm coeff}+s)\leq Z^3.
\]
The second needed comparison is \(Z\operatorname{fly}_d(k)\leq Z^2\),
using \(\operatorname{fly}_d(k)\leq Z\) and \(Z\geq1\). The price construction
chooses the atomic charge allowance so both comparisons hold.

The endpoint price-duality estimate used below has the following exact
meaning. A response duet \(P\) runs from a level \(\ell\) to \(k\leq N\)
in the same selected history. At each intermediate level its coefficient
has admissible order-zero format and remainder bound; its marked site is the
spatial origin and these sites are related by the one-step scale map. The
coefficient and cost changes satisfy the paid direction inequality, and
the densities satisfy the exact one-step transfer identity. Its left
endpoint coefficient vector is zero, and its terminal cost obeys the
spread estimate
\[
 \operatorname{cost}_k\leq
 w+c_k\bigl(w+\|P_k.x\|_{\rm coeff}\bigr),\qquad w\geq0.
\]
Write \(\kappa_k>0\) for the square root of the level-\(k\) monitor.
The endpoint estimate, which is required for all sufficiently large
period indices, uniformly in the volume index
\(N\geq k\), in the paid comparison and in \(w\geq0\), is
\begin{equation}
 \kappa_k\,\|P_k.x\|_{\rm coeff}
 \leq4(1+c_k)w.
 \label{eq:ex-trace-sharp-read}
\end{equation}
The zero in the left-end condition is a coefficient vector; the marked sites
are at the spatial origin, not zero-valued observations. The finite recurrence
and adjacent-level estimate are given in
Subsection~\ref{sec:ex-lattice-trace-price-recurrence}. The paid-cone estimate is supplied by the canonical argument in
Subsection~\ref{sec:reference-cone}, under its explicit direction, room, bank,
and endpoint-spread assumptions. None of these bounds follows from the atomic
response alone.

\subsubsection{The finite Coffer transfer}

We keep the following parameters distinct. As above, \(c_k\geq1\) is the
selected column allowance at level \(k\). Let \(X\) be the input coefficient, \(S\) its remainder budget,
\(E\) the paid response and \(e\) its attainment cost, and
\(\zeta_i>0\) its payment allowance. Let \(a\in(0,1]\) be the mask fraction,
\(v'\geq0\) the forced-allowance budget, and \(q,B>0\) two coefficient budgets.
The actual Coffer hypotheses include
\begin{equation}
 \|x_X\|_{\rm coeff}+S\leq q,\qquad
 4(1+c_k)q\leq B,\qquad e\zeta_k\leq v'B.
 \label{eq:ex-trace-coffer-inputs}
\end{equation}
They also include exact coefficient format, remainder, payment, density and
forced-profile identities; scale growth and one-step response receipts; and
the melting condition. If \(B_k\) is the spatial-spacing parameter of the
scale budget and \(R_k=\operatorname{siteFactor}(B_k)\), the latter requires
support exponent at least four and
\[
 t=\operatorname{forcedAllowance}_k
  +R_k\sqrt{\operatorname{forcedAllowance}_k}+a+R_kv',\qquad
 (2B_k+1)^4t\leq3,\qquad e^t-1\leq1.
\]
Here \(B_k\) is not the coefficient budget \(B\) in
\eqref{eq:ex-trace-coffer-inputs}.

For \(J\leq m<k<N\), let \(\operatorname{gain}(k,m)>0\) denote the
scale gain from \(m\) to \(k\). The output weight is
\begin{equation}
 w_m=a\,\frac{\zeta_m/2}
 {\|\operatorname{alto}_{m,d}\|(B/\operatorname{gain}(k,m))}.
 \label{eq:ex-trace-coffer-weight}
\end{equation}

For clarity, if \(\mathcal T_J\) is the set of reflected scale-\(J\) cells
in the volume-\(N\) tiling and \(K=|\mathcal T_J|\), with \(\mu_{N,p}\) the normalized Wilson measure, the root functional is
\[
 \operatorname{root}_{J,N,p}(F)=
 \left\|\int\prod_{x\in\mathcal T_J}
       \operatorname{copy}_x(F)(U)\,d\mu_{N,p}(U)\right\|^{1/K}.
\]
Here \(\operatorname{copy}_x(F)\) evaluates \(F\) on the translated
reflection half of the configuration; when that reflection reverses
orientation it also complex-conjugates the value. This is the
reflection-positive tiling functional, with its \(K\)th-root normalization.

Assume \(N>k\) and that the reference level is no larger than \(k\).
The response data are an admissible
sequence of coefficient coefficient families from the reference level through \(N\), an
admissible top coefficient, an exact paid response identity between each
adjacent pair, and the terminal forced-profile identity. The input field
has the format, remainder, attainment, and density identities stated above.
Assume also the melting and payment inequalities and the endpoint estimate
\eqref{eq:ex-trace-sharp-read}. Put
\[
 \widetilde G'=\operatorname{back}_{k\to N}G'.
\]
Under these hypotheses the finite construction returns one public label
\(b\), chosen before any window, with the following two properties. First,
\begin{equation}
 b'\,\|x_X-\mathcal C_k(b)\|_{\rm coeff}
   \leq 4(1+c_k)q. \label{eq:ex-trace-coffer-column}
\end{equation}
Here the window from scale \(J\) to scale \(m\) is called \emph{open} if
\[
 J\leq m,\qquad \sigma_m^2\leq A_J/A_m,
\]
where \(\sigma_m\) is the cover size of the selected clock at level \(m\)
and \(A_j\) the field dilation. Second, for the same \(b\) and for every
triple \((J,m,l)\) such that the
scale-\(J\)-to-\(m\) window is open, the reference level is no larger than
\(m\), \(m<k\), the scale \(J\) has room for the radius \(l\) at scale
\(k\), and \(G'\) is interior to the window of radius \(l\),
\begin{equation}
 \operatorname{root}_{J,N,p}(\widetilde G'-\mathcal O_{N,p}b)
   \leq \frac{\operatorname{stakes}(J,m)}{w_m}.
 \label{eq:ex-trace-coffer-root}
\end{equation}
The order of quantifiers matters: a label depending on the window would
not give a Bill, whose root estimate is required for all windows with one
label. Here \(b'\) is the endpoint price-duality constant and
\(\operatorname{root}_{J,N,p}\) is the reflection-positive root functional
on the volume-\(N\) field, defined using the scale-\(J\) tiling. In this
application the endpoint constant $b'$ is $\kappa_k$. The column
\(\mathcal C_k(b)\) is defined through the bare reference coordinate of
the public label: if \(\mathsf T\) denotes the terminal identification of
the selected response hierarchy, a linear isomorphism of the grade-\(d\)
coefficient space onto itself which sends bare reference
coordinates \(v\) to public labels, and \(I_k v\) is the coefficient of
the calibrated primary field \(\mathcal P_k(v)\) at level \(k\), then
\[
 \mathcal C_k(b)=I_k(\mathsf T^{-1}b).
\]
The construction below produces a bare coordinate \(v\) and sets
\(b=\mathsf T v\); the column estimate is proved for \(I_kv\), while the
observation \(\mathcal O_{N,p}b\) in \eqref{eq:ex-trace-coffer-root} is
that of the public label. That the level-\(k\) entries of the selected
response hierarchy, read at \(\mathsf T^{-1}b\), are these calibrated
primary fields is an identity of the selected hierarchy, used here as part
of the response data.

For the column estimate, apply the signed reference-difference construction
of Subsection~\ref{sec:reference-difference} to the Coffer coefficient field
\(X\) at terminal level \(k\), with its remainder bound \(S\) and total
charge \(q\). The construction assumes one fixed background history and
calibrated primary fields \(\mathcal P_i(v)\) whose coefficient is
\(I_i v\), together with the bank and column inequalities
\eqref{eq:reference-bank} and the exact degree-\(d\) density transition.
At each backward step it first transports the current zero-coefficient
remainder by the zero-coordinate transition
\eqref{eq:reference-zero-transition}; it then subtracts the primary field
with the resulting coefficient, so that the new remainder again has
coefficient zero. Linearity cancels the two primary densities exactly, and
the zero-coordinate transition supplies the density of the transported
remainder. The paid direction and remainder estimates
\eqref{eq:reference-direction} apply to the resulting signed difference,
not to an unsigned estimate obtained by subtracting costs.

The top family has terminal charge at most \(q\), since
\(\|x_X\|_{\rm coeff}+S\le q\). Choose its single bare reference coordinate \(v\) by
setting the left endpoint coefficient to zero, and put \(b=\mathsf Tv\).
The output is precisely
\eqref{eq:reference-produced-difference}: for the chosen coordinate, the
terminal coefficient is \(x_X-I_kv=x_X-\mathcal C_k(b)\), the terminal
density is
\(D_X-D_{\mathcal P_k(v)}\), and the left endpoint coefficient vector
vanishes. This choice does not involve any window \((J,m,l)\). Every intermediate field and the selected reference field are
evaluated against the same finest volume and background history; this is
what permits the exact density cancellation. At the terminal level the reference construction also retains its spread
bound \(r_k\le q+c_k(q+\|Z_k.x\|_{\rm coeff})\), the precise spread
condition in the Woke duet. Thus it applies \(\operatorname{SharpRead}\)
with budget \(w=q\). The stronger full terminal estimate in
\eqref{eq:reference-produced-difference}, with \(C_k=c_k\), gives
\(r_k+\|Z_k.x\|_{\rm coeff}\le4(1+c_k)q\); this is the coarse
transport budget used by the subsequent construction. The signed dualization
then gives the zero-site paid-response duet required by SharpRead, and its
endpoint inequality is exactly \eqref{eq:ex-trace-coffer-column}. Here zero refers to the left coefficient
vector and to the duet's marked spatial site being the origin. It does not
assert that the observed field values vanish.
\subparagraph{The planted polynomial and its full moment.}
Fix an admitted window $(J,m,l)$ and retain the label $b$ already chosen.
Let $V$ be its finite set of cells. Work in the square-free polynomial
algebra with generators $\mathsf t_a$, $a\in V$; for $I\subseteq V$
write $[I]Q$ for the coefficient of $\prod_{a\in I}\mathsf t_a$.
Let $D=\mathcal O_{N,p}b$ and $F=G'$, and denote their prescribed copies
at levels $N$ and $k$ by $D_a$ and $F_a$. A copy includes complex
conjugation when its reflected chamber requires it. Write $R_N$ for the
unnormalized Wilson density, $\mathsf P_{v,u}$ for density transfer
from level $u$ to $v\le u$, and $\mathsf B_{v,u}$ for its observable
adjoint. Define the following actual polynomials:
\begin{equation}\label{eq:coffer-planted-polynomial}
 Q_N=R_N\prod_{a\in V}(1-\mathsf t_aD_a),\qquad
 Q_k=\mathsf P_{k,N}Q_N,\qquad
 X_k=Q_k\prod_{a\in V}(1+\mathsf t_aF_a),\qquad
 X_m=\mathsf P_{m,k}X_k.
\end{equation}
Transfer acts coefficientwise. For bounded measurable $F$ these
coefficients are integrable: $D$ is a continuous finite-volume field,
the configuration spaces are compact, and the selected transition
kernels have bounded densities. More generally it suffices that
$\int\prod_{a\in A}|F_a|\,d\mu_k<\infty$ for every $A\subseteq V$.
Since $R_N$, the $D_a$ and the composed kernel density are bounded,
this gives absolute integrability of every mixed double integral below.

For every $I\subseteq V$, multiplication in the square-free algebra gives
\begin{equation}\label{eq:coffer-full-join}
 [I]X_k=\sum_{A\subseteq I}
  \mathsf P_{k,N}\left(R_N\prod_{a\in I\setminus A}(-D_a)\right)
                         \prod_{a\in A}F_a.
\end{equation}
Indeed a generator can occur in only one of the two factors in
\eqref{eq:coffer-planted-polynomial}; choosing the generators taken from
the second factor is exactly choosing $A$. No product of expectations
has been substituted for the transferred old product.

The room and interiority assumptions of the window put every propagated
read set inside its own chamber at each intervening step. Consequently
\eqref{eq:actual-grid-factor}, including equivariance under the prescribed
copies, gives
\[
 \mathsf B_{k,N}\Bigl(\prod_{a\in A}F_a\Bigr)
       =\prod_{a\in A}\operatorname{copy}_a
                              (\mathsf B_{k,N}F).
\]
This factorization uses disjoint kernel read sets, not just mass one.
For the integrable extension, truncate each local factor, retaining its
read set, and pass to the limit by absolute integrability and Fubini.
In the bounded case all the following interchanges follow from the
kernel-slice bounds and Fubini; in the integrable version they use the
absolute integrability just specified. Apply adjointness to each summand
of \eqref{eq:coffer-full-join} and expand the finite product. This gives
\begin{align}
 \int[I]X_k\,d\mu_k
 &=\int R_N\sum_{A\subseteq I}
       \prod_{a\in I\setminus A}(-D_a)
       \prod_{a\in A}\operatorname{copy}_a(\mathsf B_{k,N}F)
                                                    \,d\mu_N\nonumber\\
 &=Z_{N,p}\,\mathbb E_{N,p}
       \prod_{a\in I}\operatorname{copy}_a
                               (\mathsf B_{k,N}F-D).
 \label{eq:coffer-planted-moment}
\end{align}
Every further density transfer preserves the integral of every
coefficient, by its normalized mass. Therefore the left-hand side can
also be written $\int[I]X_m\,d\mu_m$. For $I=V$ this is exactly
$Z_{N,p}M$, with $M$ the moment in the root functional. For $I=\varnothing$
it is $Z_{N,p}$. This proves the full moment identity, including all
mixed terms, with the same partition function. It is stronger than a
check of the singleton densities.

\subparagraph{Constructing the two support ledgers.}
We construct the served representation on the opened centered-cell
census using the fixed label $b$ and the terminal birth in
\eqref{eq:coffer-planted-polynomial}. Throughout, the same selected
history supplies the coefficient identities, density pushes, forced
profiles, support rooms, bank and place bounds, migration, separation,
and terminal hook. The local response certificate is
\eqref{eq:coffer-response-coin}; its support and attainment bounds
are used below at their actual selected doors. The served representation
and both support ledgers will be built from these data.

Write $\tau_i$ for the weight $w_m$ marched from $m$ to $i$, and
$\Omega_i$ for the weight $w_m/a$ marched by the same positive
multipliers. Thus $\tau_i=a\Omega_i$. Put $\gamma_i=
\operatorname{forcedAllowance}_i$, and retain the column parameter
$c_i\ge1$. We use the paid scale inequalities
\eqref{eq:reference-affordable}, $\zeta_ic_i\le1$, and
$\sqrt{\gamma_i}\le1$. The latter is the small-remainder regime of
the selected history. The signed comparison constructed above has
\[
 r_i+\|Z_i.x\|_{\rm coeff}\le B/\operatorname{gain}(k,i),
 \qquad Z_k.x=x_X-I_kv.
\]
Also $\|I_kv\|_{\rm coeff}\le2B$: the triangle inequality gives
$q+B\le2B$, since $4(1+c_k)q\le B$. The two starting payments are
\begin{equation}\label{eq:coffer-double-payment}
 w_m\|\operatorname{alto}_{m,d}\|\frac{B}{\operatorname{gain}(k,m)}
       =\frac{a\zeta_m}{2}\le\frac{\zeta_m}{2},\qquad
 \frac{w_m}{a}\|\operatorname{alto}_{m,d}\|
       \frac{B}{\operatorname{gain}(k,m)}=\frac{\zeta_m}{2}.
\end{equation}
Apply the paid scale induction to both weights. It gives the
coefficient-plus-remainder payment for the signed comparison at
weight $\tau_i$, and at level $k$ gives
\[
 \Omega_k\|\operatorname{alto}_{k,d}\|B\le\zeta_k/2,
 \qquad
 \Omega_k\|\operatorname{alto}_{k,d}\|\|I_kv\|_{\rm coeff}
       \le\zeta_k.
\]
The second inequality pays the negative reference label in the
terminal birth; the label remains $-b$ for every window. The larger
weight $\Omega_k$ is essential; payment at $\tau_k$ alone would not
imply it. The first inequality pays the local response because
\[
 \tau_k\|\operatorname{alto}_{k,d}\|e
 \leq a\zeta_k e/(2B)\leq v'.
\]
Linearity of the reference hierarchy makes the bond of $-b$ equal to
that of $b$. The opened cell embeddings give occupancy one,
migration, and separation; the old-level multiplicity bound is the
admitted schedule.

Apply the terminal birth construction and its finite downward
induction to this negative label, starting with the second payment.
The terminal birth expansion starts with the actual background and
linear insertions and transports them through the selected transition
identities, retaining the forced hook. At the finest level its
polynomial is precisely $Q_N$ in
\eqref{eq:coffer-planted-polynomial}. The construction begins with the background box factors and linear
terminal insertions, moves forced responses into boxes by connected
component closure, and removes the constant and singleton
coefficients. This defines the higher arrays, rather than merely
bounding proposed arrays. More explicitly, the terminal place bound has a footprint of cover
number at most one and section budget $t_N\ge1$. The bank inequality is
\[
 e^{\alpha_N+2}\mathfrak c_N(t_N)\zeta_N\le u_N.
\]
Together with the paid negative-label bond this gives each terminal
response its weighted price $u_N$ and its prescribed forced hook.
Writing $n_N$ for the terminal occupancy bound, the terminal budget
requires
\[
 C_N=\gamma_N+n_N\operatorname{siteFactor}(B_N)u_N,
 \quad (2B_N+1)^4C_N\le3,
 \quad e^{C_N}-1\le1,\quad \alpha_N\ge3.
\]
The place clause here is the cover-number-at-most-one receipt. The
birth polynomial is integrable by the compact-continuity estimate for
the terminal cylinder observable. The mixed polynomial is integrable under each transition: on this
finite-volume compact configuration space its cylinder factors are
bounded, and the selected transition kernels are normalized with
bounded densities.
Thus the component construction of \eqref{eq:transfer-closure} and
\eqref{eq:ope-forest-price}, with spare exponent three, gives a flagged
higher array of price at most one at exponent $\alpha_N-3$.
The ordinary exponent started linear and remains linear after the
response shift, so its higher array is zero. The terminal hook and
singleton coin identities identify the prescribed singletons; the
higher-degree projection retains the exact birth polynomial.
At level $k$, the seeded negative polynomial $Q_k$ has the negative
reference singletons and higher arrays of price one at weights
$\Omega_k$. Reweighting only degrees at least two from $\Omega_k$
to $\tau_k=a\Omega_k$ reduces those prices to $a$. Multiply by the
positive factors $\prod_a(1+\mathsf t_aF_a)$ and insert their local
response, whose singleton coin and forced profile are
\eqref{eq:coffer-response-coin}. This forms the mixed polynomial
$X_k$, with singleton equal to positive primary plus negative reference
plus response. Before component closure, the ordinary higher array
still has price at most $a$; the flagged background, forced profiles
and reweighted old boxes are priced explicitly below.
For each nonzero term, connect its old flagged boxes and response
supports by support intersection, require an old flagged box, and put
the product at the union support with the union flag set. This finite
component regrouping preserves the full polynomial. A component with
one response contributes exactly its forced hook; components with two
or more responses have degree at least two. The resulting representation
of $X_k$ has ordinary and flagged higher arrays $H^0_Y,K^0_P$, each
with site price at most one at exponent $\alpha_k-4$. The flagged
price is the component bound proved below; the ordinary higher array
remains at price $a$.

Here and below a site price $C$ means that the weighted absolute
coefficient sums have nonnegative majorants $h_Y$ satisfying
$\sum_{Y\ni x}h_Ye^{\alpha c_{B}(Y)}\le C$ for every $x$;
for boxes use their first, support component. This includes a
majorant uniform in the configuration. All higher arrays have
marker degree at least two. Changing their weights from
$\Omega_k$ to $\tau_k=a\Omega_k$ multiplies the price of a degree
$j\ge2$ coefficient by $a^j\le a$. Consequently both old higher
arrays now have site price at most $a$.

Let $\phi_a$ and $q_a$ denote the negative reference ordinary
singleton and forced profile, respectively. The actual atomic
response, after the prescribed motion and dilation, gives an
ordinary field $h_a$, a dominated forced family $y_a$, and a
support response $\psi_{a,Y}$, with
\begin{equation}\label{eq:coffer-response-coin}
 \operatorname{coin}_k(h_a,y_a)=\mathscr R_k F_a,
 \qquad \operatorname{profile}_k(y_a)=\mathcal H_k(\psi_a).
\end{equation}
Here $\mathscr R_k$ denotes the evaluation of the level-$k$ representation,
and $\mathcal H_k$ is the response hook defined in
Subsection~\ref{sec:local-transfer}. These identities follow by applying
the same motion and dilation to the input density and forced-profile
identities: the Wilson density is invariant under that motion, and
both profile and hook are linear. The response has a majorant
$g_{a,Y}$ with
\[
 \sum_Y\tau_k g_{a,Y}e^{\alpha_k c_{B_k}(Y)}\le v'.
\]
Indeed the undilated attainment cost is $e$, and
$\tau_k\|\operatorname{alto}_{k,d}\|e\le
 a\zeta_k e/(2B)\le v'$, by
\eqref{eq:coffer-double-payment} and $e\zeta_k\le v'B$.
The attainment exponent $\alpha_k+2$ is stronger than the exponent
used here. Each nonzero response contains its mark and vanishes
when its support is healthy. The census marks at $k$ are distinct
by the admitted window separation, so occupancy there is one.
Equation~\eqref{eq:transfer-root-hit}, with exponent $\alpha_k-1$,
therefore bounds the combined response site price by $R_k^{\rm site}v'$,
where $R_k^{\rm site}=\operatorname{siteFactor}(B_k)$.

To avoid confusing this packing factor with the density, write
$R^{\rm site}=R_k^{\rm site}$ for the rest of the gluing step.
The forced profiles $q_a$ of the negative reference field have
site price at most $R^{\rm site}\sqrt{\gamma_k}$ at exponent
$\alpha_k-1$. Here is the payment check. The reference remainder
has bound at most $c_k\|I_kv\|_{\rm coeff}$; its forced family
has cost at most $\|\operatorname{alto}_{k,d}\|\sqrt{\gamma_k}$
times that remainder. Multiplication by $\tau_k$, the reference
bond payment, and $\zeta_kc_k\le1$ give the per-mark cost
$\sqrt{\gamma_k}$. Apply \eqref{eq:transfer-root-hit} again.
The background box activity has site price $\gamma_k$. Thus the
full old box array, including its background, its linear forced
profiles and $K^0$, has site price at most
\[
 W=\gamma_k+R^{\rm site}\sqrt{\gamma_k}+a
                       \quad\hbox{at exponent }\alpha_k-1.
\]

Multiplication by $\prod_a(1+\mathsf t_aF_a)$ changes the old
ordinary exponent by $\sum_a\mathsf t_aF_a$. Put
$J=\sum_{a,Y}\mathsf t_a\psi_{a,Y}$. Form a new box array
$\overline K$ as follows. Take finite sets of distinct old box
addresses and response addresses $(a,Y)$, with connected support
intersection graph, at least one old box, and pairwise disjoint
old box supports. Assign their product to the box whose support
is their union and whose flag set is the union of the old flag
sets, retaining the indicator that this is the flag set of the
union support. Sum the products with the same output box. This is
an explicit finite definition on each finite-volume lattice.
The component decomposition in \eqref{eq:transfer-closure} gives
\[
 \operatorname{field}(\overline K)
   =\operatorname{field}(K^{\rm old})\operatorname{pass}(J).
\]
Subtract $J$ from the ordinary exponent. The represented polynomial
is still exactly $X_k$; no coefficient has been discarded.

Set $d_*=\alpha_k-4\ge0$ and
$\eta=W+R^{\rm site}v'$. The melting hypotheses say
$(2B_k+1)^4\eta\le3$ and $e^\eta-1\le1$.
Apply \eqref{eq:ope-forest-price} with output exponent $d_*$
and spare exponent $3$. Its incoming exponent is
$d_*+3=\alpha_k-1$, exactly the exponent just paid. The price
of a product is at most the product of its coefficient majorants,
and the union box is uniquely determined by its constituent
addresses. Restricting the connected-set sum to the admissible
sets above only decreases its nonnegative majorant. Hence
$\overline K$ has site price at most $e^\eta-1\le1$ at exponent
$d_*$. The ordinary higher array remains $H^0$, since $J$ and
all added ordinary terms are linear; its price is at most
$a\le1$ at the same exponent.

The vacuum coefficient of $\overline K$ is the old background.
Its singleton at mark $a$ is $q_a+\mathcal H_k(\psi_a)$:
without a response a component is one old box, while a component
with exactly one response gives precisely the hook, as proved in
Subsection~\ref{sec:local-transfer}. Equation
\eqref{eq:coffer-response-coin}, the singleton identity for $Q_k$,
and the product rule in the square-free algebra identify the
singleton of $X_k$ with the coin of $(\phi_a+h_a,q_a+
\operatorname{profile}_k(y_a))$. Positivity of the background
density and the singleton uniqueness in
\eqref{eq:ope-single-coin} then identify the ordinary singleton
of the new exponent. Remove degrees zero and one from each
array; the remaining $H^k,K^k$ give exactly this prescribed
singleton representation of $X_k$. Degree projection cannot
increase the absolute coefficient price. Products, finite sums
and local flag indicators preserve measurability, gauge
invariance, locality and containment of each occurring mark in
its support. Both higher ledgers therefore have price one at
exponent $\alpha_k-4$.

Now descend from $k$ to $m$ along the same selected history. At each
step the signed singleton expansion in Subsection~\ref{sec:reference-difference}
has its exact density push, forced-profile identity and finite support
witness on the chosen door. Transporting the higher arrays through that
witness preserves the full polynomial; the newly added response terms
have degree at least two, and their connected supports are grouped by
union. Thus the descent transports the representation already
constructed at $k$ rather than postulating a new served object.
For specificity the required paid local data are
\[
 0<\delta_i\le1,\quad \alpha_i,\alpha_{i+1}\ge100,\quad
 T_i=n_i\operatorname{siteFactor}(B_{i+1})\delta_i+4\delta_i,
 \quad T_i+e^{T_i}T_i\le M_i,
\]
\[
 A_i=u_i+n_i\operatorname{siteFactor}(B_i)\delta_i,
 \qquad (2B_i+1)^4A_i\le\tfrac12,
 \qquad e^{A_i}-1\le1.
\]
Here $M_i$ is the available input budget, $u_i$ is the flagged
response price, and $n_i$ bounds the two adjacent occupancies. These
are the one-step payment inequalities of the selected history, not
additional conclusions of melting at $k$.
Here is the singleton payment check. Put
$C_i^{\rm wing}=\operatorname{orb}(d)(\psi_i+2c_{i+1})\ge2$.
The retained signed expansion has dilated fine cost
\[
 s_i\le\|\operatorname{alto}_{i+1,d}\|
       \bigl(C_i^{\rm wing}\|Z_{i+1}.x\|_{\rm coeff}
                    +2r_{i+1}\bigr)
 \le \|\operatorname{alto}_{i+1,d}\|C_i^{\rm wing}
          (\|Z_{i+1}.x\|_{\rm coeff}+r_{i+1}).
\]
The paid orbit inequality of the selected history is
$2C_i^{\rm wing}\zeta_{i+1}\le1$; hence
$2\tau_{i+1}s_i\le1$. Its coarse response cost is at most
$t_i\le\|\operatorname{alto}_{i,d}\|r_i$. Since the march has
$\tau_{i+1}\le\tau_i$ and $\zeta_i\le1$, the same paid comparison
gives $\tau_{i+1}t_i\le1$. These pay the input and output sides of the actual finite-support
comparison. The parent-site identity, occupancy and cover-price
comparison are the same migration and separation data as in the
centered induction. The adjacent transition preserves the singleton
entries; every added response term has marker degree at least two.
The component-closure estimate therefore recovers exponent $\alpha_i$
and separately bounds both higher ledgers by one. At the last level, the forced remainder allowance
$\sqrt{\gamma_m}$ and the marched coefficient-plus-remainder
payment give the Vault constants $2,1$. The first step accepts
$\alpha_k-4$; every subsequent step accepts the full exponent
already obtained. The actual polynomial at the next level is
$\mathsf P_{i,i+1}X_{i+1}$, so finite induction constructs the
required support families for the precise $X_m$ in
\eqref{eq:coffer-planted-polynomial}.

Finally the singleton bounds are separate from the higher-array
bounds. At $m$, let $z_m=\|Z_m.x\|_{\rm coeff}$. The coefficient
format and remainder bounds imply, before dilation and motion,
\[
 |\operatorname{exponentialEntry}(Z_m)|\le(1+c_m)z_m+r_m,
 \qquad \sum_P|\operatorname{profile}(Z_m^{\rm f})_P|\le r_m.
\]
For the first inequality, sum the section bounds against the
absolute coefficient vector and add the remainder bound; the
constant coordinate contributes the extra $z_m$. The second
is the absolute forced-profile bound, since
$\sqrt{\gamma_m}\le1$. Motion preserves absolute values.
Since $w_m\|\operatorname{alto}_{m,d}\|(r_m+z_m)\le\zeta_m$
and $\zeta_mc_m\le1$, the ordinary singleton has weighted
absolute value at most $2$ and the summed flagged singleton at
most $1$. These are exactly the two singleton bounds used by the finite-census
price estimate below. Together with the constructed higher arrays and
their site budgets, they give the required served representation and
its support bounds at level $m$.

Thus the following support decomposition has been constructed by the
paid same-history transition and its finite-support witnesses.
Let \(V\) be the set of census sites, \(K=|V|\), with marks \(o_a\),
\(a\in V\), in the level-\(m\) lattice. The polynomial $\mathsf X=X_m$ of
\eqref{eq:coffer-planted-polynomial}, whose integrability and moment
identity were just established, has the following structure. There are singleton site
amplitudes \(\Phi_a(U)\) and singleton box amplitudes \(\omega_{a,P}(U)\),
\(a\in V\), \(P\) a box, with
\[
 w_a|\Phi_a(U)|\leq2,\qquad \sum_P w_a|\omega_{a,P}(U)|\leq1
 \qquad\text{for every }a\text{ and every }U,
\]
and two correction families, an ordinary family \(H_Y(U)\) indexed by
supports \(Y\) and a flagged family \(K_P(U)\) indexed by boxes \(P\), such
that
\[
 \mathsf X(U)=e^{E_m(U)}\,
   \operatorname{pass}\Bigl(\sum_a\Phi_a(U)\,\mathsf t_a+\sum_YH_Y(U)\Bigr)
   \cdot\operatorname{field}\Bigl(P\mapsto
       \mathsf b_P(U)+\sum_a\omega_{a,P}(U)\,\mathsf t_a+K_P(U)\Bigr),
\]
where \(E_m\) is the exponent of the level-\(m\) effective density,
\(\mathsf t_a\) is the census monomial of the site \(a\), \(\mathsf b_P\)
is the box activity term of the level-\(m\) representation, and
\(\operatorname{field}\) is evaluated on the flagged set of \(U\). Each
correction family satisfies the anchored support-cost bound with constant
\(1\) at the support exponent \(\alpha_m\) of the level-\(m\) budget: there
are numbers \(h_Y\geq0\) bounding the weighted absolute coefficient sums,
\(\mathcal P_w(H_Y(U))\leq h_Y\) for all \(U\), with
\(\sum_{Y\ni x}h_Ye^{\alpha_m\operatorname{cov}_m(Y)}\leq1\) for every
site \(x\), and likewise for \(K_P\) with the support of the box \(P\).
Every nonzero correction term has a census mark \(o_a\) in its support.
Consequently, summing the anchored bound over the \(K\) marks,
\begin{equation}
 \sum_Y\mathcal P_w(H_Y(U))\leq K,\qquad
 \sum_P\mathcal P_w(K_P(U))\leq K
 \qquad\text{for every }U.
 \label{eq:ex-trace-coffer-corrections}
\end{equation}
The required support-price structure is a statement about this precise
polynomial $X_m$. Its full moment identity is
\eqref{eq:coffer-planted-moment}, already proved independently of those
price bounds. The support decomposition and its two prices were constructed above,
using the explicitly stated local-transition inputs.

The polynomial-price calculation behind the finite-gas estimate is as
follows. For a finite census polynomial \(Q\), write
\[
 \mathcal P_w(Q)=\sum_s\operatorname{toll}(w,s)\|Q_s\|
\]
for its weighted absolute coefficient sum. The price is submultiplicative.
For box terms \(K_p\), the finite sum over disjoint-support patterns obeys
\[
 \mathcal P_w(\operatorname{field}(K))
 \leq\prod_p(1+\mathcal P_w(K_p))
 \leq\exp\!\left(\sum_p\mathcal P_w(K_p)\right),
\]
because each pattern's product price is at most the product of its factors
and \(1+x\leq e^x\). The same submultiplicativity, applied term by term to
the finite truncated exponential, gives
\(\mathcal P_w(\operatorname{pass}(Q))\leq e^{\mathcal P_w(Q)}\).
The singleton site and box amplitudes contribute \(2K\) and \(K\) to the
exponent, by the two singleton bounds just proved. The
ordinary and the flagged correction families contribute \(K\) each, by
\eqref{eq:ex-trace-coffer-corrections}. The total is \(5K\); the box
activity terms \(\mathsf b_P\) and the exponent \(E_m\) are paid by the
finite-volume denominator bound used next.

By \eqref{eq:coffer-planted-moment}, the integral of the selected coefficient
equals the volume-$N$ Wilson partition function times $M$.
To specify the denominator estimate, let $A_m$ be the lawful local
representative used in this very expansion, and take
$\rho=\widehat\beta_m^{-1/2}\le\min\{1,\vartheta_m\}$.
The numerator estimate just obtained is the coefficient-price estimate
of Subsection~\ref{ex-centered-construction}, with total mark charge $5K$.
Its small-ball lower bound for $Z_{A_m}=\int\operatorname{eval}(A_m)$,
followed by division as in \eqref{ex-centered-sail}, bounds the quotient
by
\[
 \exp\!\left(|\Lambda_m|\left[2\operatorname{exCost}_m+
 \kappa_m+\beta_m\operatorname{baseBudget}_m(\rho)
 +4d_G\log(16/\rho)\right]+5K\right).
\]
Here $\kappa_m$ denotes the vacuum price in that estimate, not the
coefficient-recovery constant used later. The treasury receipt at this
same radius bounds the bracket by $\operatorname{Treasury}(m)$.
The same-history transition-mass identity identifies
$Z_{A_m}=Z_{N,p}>0$, as in
Lemma~\ref{lem:ex-selected-birth-denominator}; this is also the partition
function in the planted moment identity above. No comparison of
different local representatives or partition functions is needed.
Consequently
\begin{equation}
 |M|\,\operatorname{toll}(w,\mathfrak s)
 \leq \exp\!\left(|\Lambda_m|\operatorname{Treasury}(m)
                         +K(2+1+2)\right).
 \label{eq:ex-trace-spiker}
\end{equation}
This is an absolute-value estimate and does not use cancellation. Its hypotheses include the paid local-transition inputs used to construct
the support ledgers above.

In the Coffer application all site weights equal \(w_m\), so
\(\operatorname{toll}(w,\mathfrak s)=w_m^K\). The tiling identity is
\[
 |\Lambda_m|=K\left(\frac{A_J}{A_m}\right)^4.
\]
Taking the positive \(K\)-th root of \eqref{eq:ex-trace-spiker} yields
\[
 |M|^{1/K}\leq
 \frac{\exp((A_J/A_m)^4\operatorname{Treasury}(m)+5)}{w_m}
 =\frac{\operatorname{stakes}(J,m)}{w_m},
\]
which is \eqref{eq:ex-trace-coffer-root}. The moment identity is needed to
identify this \(M\) with the desired rooted observable.

\subsubsection{Using the output in the Bill}

The Bill is evaluated on the inflated field \(G'=u g\). Its root error is
\(\widetilde G'-\mathcal O_{N,p}b\), where
\(\widetilde G'=\operatorname{back}_{k\to N}G'\). In the Bill specialization, \(a=Z^{-1}\), \(v'=Z^{-1}\), \(B=Z^3\), and \(e=Z\), while \(\zeta_i=\operatorname{fly}_d(i)\) at each scale. Thus the payment condition is \(e\zeta_k=Z\operatorname{fly}_d(k)\leq Z^2=v'B\). The gas weight
from \eqref{eq:ex-trace-coffer-weight} is
\[
 w_{\rm gas}=Z^{-1}
 \frac{\operatorname{fly}_d(m)/2}
 {\|\operatorname{alto}_{m,d}\|(Z^3/\operatorname{gain}(k,m))}.
\]
The factor $Z^{-1}$ is the mask fraction $a$ in the Coffer input, so this is
exactly the native Bill denominator.

The melting inequalities also hold eventually for this specialization.
To see this directly, write $x=\widehat\beta_k$, $B_x=\lceil(\log x)^2\rceil$
and $R_x=\exp(1)[2(16B_x+9)]^4$ for the spatial radius and site factor.
The selected forced allowance is exactly
\[
 f(x)=\exp\{-\exp[(\log\log x)^2/5]\}
\]
for sufficiently large $x$. Indeed, the monitor is
$H(x)=\exp((\log\log x)^2)$ and the allowance is
$\exp(-H(x)^{1/5})$; the power $1/5$ is applied to the value of the
monitor, not to its argument. With $y=\log\log x$,
$\exp(y^2/5)/\log x=\exp(y^2/5-y)\to\infty$.
Consequently, for every $A>0$, $\sqrt{f(x)}\le x^{-A}$ eventually;
also $0<f(x)\le\sqrt{f(x)}$ eventually. With $a=v'=x^{-1}$, the
quantity in the melting condition is
\[
 t_x=f(x)+R_x\sqrt{f(x)}+x^{-1}+R_xx^{-1}.
\]
Here $R_x=O((\log x)^8)$ and $(2B_x+1)^4=O((\log x)^8)$, so
$(2B_x+1)^4t_x\to0$ and $t_x\to0$.
Thus eventually $(2B_x+1)^4t_x\le3$ and $e^{t_x}-1\le1$;
$0<a\le1$ and $v'\ge0$ hold as well. The remaining clause is the
support exponent bound $\alpha_k\ge4$, retained among the scale
hypotheses. All these tails depend only on the fixed scale data, before
periods, volumes, atomic certificates and windows are chosen.

The other two entries of the Bill use the specific atomic charge. Fix
\(u>0\) and let the charge in \eqref{eq:ex-trace-atomd-charge} be
\[
 q=q_k=u\,\vartheta_k^2\operatorname{Age}_k(1/6),
 \qquad
 \operatorname{Age}_k(a)=\exp\bigl(a(\log\log\widehat\beta_k)^2\bigr),
\]
where \(\vartheta_k\) is the threshold of the scale-\(k\) budget; thus
\(\kappa_k=\operatorname{Age}_k(1/2)\) and
\(\operatorname{Age}_k(a)\operatorname{Age}_k(a')=\operatorname{Age}_k(a+a')\).
The selected schedule has \(c_k\leq\log\widehat\beta_k\). Since
\(\vartheta_k^2\operatorname{Age}_k(1/6)=\widehat\beta_k^{-1}
\operatorname{Age}_k(2+1/6)\) tends to zero faster than any power of
\(\log\widehat\beta_k\) grows, and
\(\log\widehat\beta_k=\exp(\log\log\widehat\beta_k)\) is eventually
smaller than any positive multiple of \(\operatorname{Age}_k(1/12)\), and
\(\widehat\beta_k\to\infty\), for all large \(k\)
\[
 q_k\leq1,\qquad 4(1+\log\widehat\beta_k)q_k\leq Z^3,\qquad
 4(1+\log\widehat\beta_k)u\leq\operatorname{Age}_k(1/12);
\]
the second inequality gives \eqref{eq:ex-trace-coffer-budget}. The norm
entry of the Bill is \(\|x_X\|_{\rm coeff}\leq1\); it follows from
\eqref{eq:ex-trace-atomd-charge}, nonnegativity of the remainder allowance
and \(q_k\leq1\):
\(\|x_X\|_{\rm coeff}\leq\|x_X\|_{\rm coeff}+s\leq q_k\leq1\). The bias
entry is obtained from \eqref{eq:ex-trace-coffer-column} with
\(b'=\kappa_k\) alone:
\[
 \|x_X-\mathcal C_k(b)\|_{\rm coeff}
 \leq\frac{4(1+c_k)q_k}{\kappa_k}
 \leq\operatorname{Age}_k(1/12)\,\vartheta_k^2
      \operatorname{Age}_k(1/6)\operatorname{Age}_k(-1/2)
 =\vartheta_k^2\operatorname{Age}_k(-1/4),
\]
because \(1/12+1/6-1/2=-1/4\). This is the bias bound which
Subsection~\ref{sec:ex-lattice-trace-bill-prices} takes as a hypothesis on a
paid label; that section does not prove it.

The resulting statement has the following order of
choices, under the standing hypotheses on the selected family, which
include the served response hierarchy and its terminal identification,
the same-history finite-menu and atomic response receipts, eventual
support exponents at least four, the column bound
$c_k\le\log\widehat\beta_k$, and the endpoint estimate
\eqref{eq:ex-trace-sharp-read} with
\(b'=\kappa_k\) as a hypothesis. For every fixed \(u>0\) and grade
\(d\): for all sufficiently
large \(k\), for all sufficiently large periods \(p\) (depending on \(k\)),
for every volume index \(N>k\) and every atomic certificate \((X,g)\) at
scale \(k\) with charge \(q_k\) and attainment cost at most
\(Z=\widehat\beta_k\), there is a public label \(b\) such that
\((g,x_X,b)\) is a Bill, with the root estimate for all windows. The
period threshold is the one at which the endpoint estimate
\eqref{eq:ex-trace-sharp-read} holds with \(\ell\) the reference level.

The remaining analytic dependencies are the atomic response, the endpoint
price-duality recurrence, and the paid elementary integration and singleton
transition inputs specified above. The site/box expansion, its two separate
correction prices, and its singleton envelope are now constructed from
those inputs; no additional served polynomial is postulated.
